\chapter{Phân tích mạng và kết quả}
\label{ch:analysis}
Phân tích tách năm phép chiếu: đồng tác giả, hợp tác tổ chức, trích dẫn có hướng
trong tập dữ liệu, bài báo--chủ đề và đóng góp tổ chức--chủ đề. Các chỉ số lấy
từ phân tích đồ thị, tính ngoại tuyến bằng NetworkX \citep{networkxDocs}.
Mỗi phép chiếu giữ các nút thuộc phạm vi, kể cả nút cô lập; mạng hợp tác là vô
hướng, còn cạnh trích dẫn đi từ bài trích dẫn tới bài được trích dẫn.

\section{Xu hướng công bố và chủ đề nghiên cứu}
Số liệu chính thức ở Hình~\ref{fig:publications} cho thấy năm 2025 có nhiều
bài hơn năm 2021 khoảng 80\%, với bước giảm vào năm 2022. Sự mở rộng không
đều và chỉ phản ánh các tập kỷ yếu luồng chính đã chọn.

Góc nhìn chủ đề trả lời câu hỏi khác: phần có nhãn của tập dữ liệu phân bố thế
nào giữa các lĩnh vực do mô hình định nghĩa? Bài có $L$ nhãn được chấp nhận
đóng góp $1/L$ cho mỗi nhãn. Tỷ trọng hằng năm chia tổng phần đóng góp cho số
bài có nhãn, nên một bài đa nhãn không bị tính thành nhiều bài trọn vẹn.
Hình~\ref{fig:topics} gộp hai hội nghị và chọn tám chủ đề có tổng phần đóng góp
lớn nhất.

\plot{topic-trends}{Xu hướng suy ra của tám chủ đề có tổng phần đóng góp lớn
nhất. Mỗi bài có nhãn đóng góp tổng cộng một đơn vị, chia giữa các nhãn; mẫu số
hằng năm là tổng bài có nhãn AAAI và IJCAI. Các ô dùng chung thang trục đứng.}{fig:topics}
\FloatBarrier

Học biểu diễn nổi bật trong toàn giai đoạn, còn tỷ trọng học đa phương thức
suy ra tăng ở cả hai hội nghị. Trong bảng riêng từng hội nghị, tỷ trọng đa
phương thức của AAAI tăng từ 6.96\% năm 2021 lên 16.42\% năm 2025; IJCAI tăng
từ 7.43\% lên 12.90\%. Cùng chiều thay đổi ở hai chuỗi hội nghị khiến đây là
một tín hiệu đáng khảo sát, bên cạnh tác động do cơ cấu hội nghị thay đổi.
Tuy nhiên, kết quả vẫn là biến động tỷ trọng nhãn mô hình trong nhóm bài được
gán nhãn. Các bài bị từ chối gán nhãn và giới hạn đánh giá ở Chương~\ref{ch:quality}
hạn chế kết luận về toàn bộ bức tranh nghiên cứu.

\section{Mạng hợp tác giữa các nhà nghiên cứu}
Phép chiếu đồng tác giả nối \Researchers{} nhà nghiên cứu bằng
\CoauthorPairs{} cặp phân biệt. Bậc đếm số người cộng tác khác nhau, còn độ
mạnh phân bổ cộng dồn hỗ trợ sau điều chỉnh theo quy mô nhóm tác giả. Bậc trung
bình là 5.00, lớn nhất là 90. Giá trị trung bình thấp đi cùng đuôi trên dài cho
thấy một nhóm nhỏ định danh có phạm vi cộng tác rộng hơn đáng kể so với phần
đông nhà nghiên cứu.

\plot{researcher-structure}{Cấu trúc mạng đồng tác giả. Trái: tỷ lệ trong toàn
bộ 32,189 nhà nghiên cứu có bậc ít nhất $k$; chỉ vẽ bậc dương trên trục log,
nhưng mẫu số vẫn giữ nút cô lập. Phải: toàn bộ 2,936 thành phần xếp theo kích
thước, gồm cả nút cô lập; trục đứng dùng thang log. Đường cong không đủ xác
lập mô hình luật lũy thừa.}{fig:structure}

Thành phần liên thông lớn nhất gồm \LargestComponent{} nhà nghiên cứu, tương
đương 64.5\% phép chiếu. Tổng cộng có \ResearcherComponents{} thành phần và
\ResearcherIsolates{} nhà nghiên cứu cô lập. Đường đi trong thành phần lớn
nhất thể hiện một chuỗi quan hệ đồng tác giả, không có nghĩa mọi thành viên
đều trực tiếp làm việc cùng nhau. Một thành phần lớn bên cạnh nhiều nhóm nhỏ
cho thấy lõi kết nối rộng cùng sự phân mảnh cục bộ đáng kể.

Chiều thời gian bổ sung ngữ cảnh. Bảng~\ref{tab:annual} đếm nhà nghiên cứu có
ít nhất một quan hệ tác giả được chấp nhận trong mỗi năm, đồng thời gán từng
cặp đồng tác giả về năm sớm nhất trong các bài hỗ trợ. Một người có thể xuất
hiện ở nhiều dòng. Cột cuối đo cặp được quan sát lần đầu trong cửa sổ năm năm,
không phải mọi cặp hoạt động trong năm hay lần đầu cộng tác trong sự nghiệp.

% Generated from frozen v4; edit generator.
\begin{table}[!htbp]\centering\small
\caption{\lang{Annual activity in the accepted coauthorship network.}{Hoạt động hằng năm trong mạng đồng tác giả được chấp nhận.}}\label{tab:annual}
\begin{tabularx}{\textwidth}{rYY}
\toprule \lang{Year}{Năm} & \lang{Active researchers}{Nhà nghiên cứu hoạt động} & \lang{First-observed pairs}{Cặp xuất hiện lần đầu}\\\midrule
2021 & \num{6337} & \num{12599}\\
2022 & \num{6136} & \num{10730}\\
2023 & \num{6941} & \num{12998}\\
2024 & \num{9361} & \num{17447}\\
2025 & \num{13082} & \num{26708}\\
\bottomrule\end{tabularx}\end{table}

\FloatBarrier

Số định danh hoạt động tăng từ 6,337 năm 2021 lên 13,082 năm 2025; số cặp xuất
hiện lần đầu tăng từ 12,599 lên 26,708. Kết hợp với xu hướng công bố, các số
liệu cho thấy phần cuối giai đoạn có nhiều bài hơn và mạng hợp tác được chấp
nhận rộng hơn. Việc giữ lại định danh có thể loại cạnh không đồng đều, nên đồ
thị không tách được hoàn toàn biến động hoạt động thực tế khỏi ảnh hưởng của
độ bao phủ siêu dữ liệu.

\section{Mạng hợp tác giữa các tổ chức}
Phép chiếu tổ chức gồm \Organizations{} tổ chức và \OrganizationPairs{} cặp
hợp tác duy nhất. Có \OrganizationIsolates{} nút cô lập, bậc trung bình 12.73
và bậc lớn nhất 257. Mỗi cạnh cần bằng chứng đơn vị công tác theo cùng bài của
hai tác giả khác nhau tới hai tổ chức khác nhau. Như vậy, quan hệ cùng công bố
được xác định bằng quy tắc cụ thể.

\textbf{Cặp duy nhất khác với lượt hỗ trợ cặp--bài.} Một cặp tổ chức chỉ tạo
một cạnh dù xuất hiện trong nhiều bài. Số bài hỗ trợ đếm các bài phân biệt.
Một bài có nhiều cặp đủ điều kiện đóng góp một lượt hỗ trợ cho mỗi cặp; trọng
số phân bổ còn chia tổng đóng góp của bài giữa các cặp. Ma trận khu vực ở
Hình~\ref{fig:sectors} cộng lượt hỗ trợ cặp--bài, không đếm bài duy nhất hay
cạnh duy nhất. Hai ô đối xứng biểu diễn cùng lượt hỗ trợ liên khu vực, không
được cộng hai lần.

\plot{sector-mixing}{Hợp tác tổ chức theo khu vực hoạt động. Mỗi ô đếm lượt hỗ
trợ cặp tổ chức--bài thuộc tập lõi. Ma trận đối xứng giữ nhóm chưa rõ khu vực;
mỗi cặp đủ điều kiện cần hai tác giả khác nhau có bằng chứng tổ chức tương ứng
trên cùng bài.}{fig:sectors}
\FloatBarrier

Học thuật--học thuật là ô lớn nhất, với 10,956 lượt hỗ trợ. Học thuật--doanh
nghiệp có 3,533 lượt. Kết quả cho thấy cùng công bố giữa các tổ chức học thuật
là kiểu phổ biến nhất trong dữ liệu, đồng thời kết nối liên khu vực cũng hiện
diện đáng kể. Khu vực là thuộc tính tổ chức do OpenAlex cung cấp; siêu dữ liệu
chưa rõ được giữ tường minh. Số liệu thể hiện hỗ trợ công bố, không xác lập
quan hệ hợp đồng.

\Needspace{17cm}
Hình~\ref{fig:orgpairs} chọn mười hai cặp học thuật--doanh nghiệp có nhiều bài
hỗ trợ phân biệt nhất trên mọi chủ đề. Cách trình bày này chuyển mạng dày thành
so sánh dễ đọc về cùng công bố lặp lại. Mười hai cặp gồm mười một tổ chức duy
nhất, nên một tổ chức có thể xuất hiện ở nhiều dòng. Đây không phải đồ thị con
cảm sinh trên các tổ chức đó.

\plot{organizations-network}{Mười hai cặp học thuật--doanh nghiệp dẫn đầu theo
số bài hỗ trợ phân biệt được chấp nhận, gộp mọi chủ đề. Trường hợp bằng nhau
được xếp theo ID hai đầu. Tổ chức lặp giữa các dòng là có chủ ý; mỗi dòng là
một cặp duy nhất.}{fig:orgpairs}
\FloatBarrier

Góc nhìn theo chủ đề cung cấp ví dụ cụ thể hơn. Shanghai Jiao Tong
University--Tencent có 14 bài hỗ trợ phân biệt mang nhãn \emph{Học biểu diễn},
đồng hạng với Xiamen University--Tencent và Zhejiang University--Alibaba. Hai
bài của cặp Shanghai Jiao Tong--Tencent đã được kiểm tra nguồn là
\href{https://doi.org/10.1609/aaai.v36i1.19954}{\emph{Exploiting Fine-Grained
Face Forgery Clues via Progressive Enhancement Learning}} và
\href{https://doi.org/10.1609/aaai.v36i3.20193}{\emph{DIRL: Domain-Invariant
Representation Learning for Generalizable Semantic Segmentation}} (2022).
Các quan sát được chấp nhận nối hai tổ chức với hai tác giả khác nhau.
Kết quả 14 bài áp dụng bộ lọc chủ đề, còn Hình~\ref{fig:orgpairs} gộp mọi chủ đề.

Hồ sơ chủ đề của tổ chức cũng phụ thuộc mẫu số. Với chủ đề suy ra Học biểu
diễn, Zhejiang University có 156 bài hỗ trợ trong 489 bài có nhãn gắn với tổ
chức; Peking University có 139 trong 381 bài. Tỷ trọng hồ sơ theo phân bổ lần
lượt là 25.1\% và 28.8\%. Mỗi bài có nhãn được chia đóng góp giữa các tổ chức
hỗ trợ và giữa các nhãn, nên các tỷ trọng khác với tỷ số thô 156/489 và
139/381. Trong phần hồ sơ quan sát được, Zhejiang có nhiều bài hỗ trợ hơn,
còn Peking dành tỷ trọng phân bổ lớn hơn cho chủ đề này. Đây là so sánh hữu
ích giữa quy mô và mức tập trung, với điều kiện dựa trên liên kết tổ chức được
giữ lại và nhãn suy ra.

\section{Cấu trúc cộng đồng và độ trung tâm của nhà nghiên cứu}
Phát hiện cộng đồng tìm các nhóm có kết nối đồng tác giả tương đối dày bên
trong. Lần chạy Louvain được chọn dùng hạt giống 17 và độ phân giải 0.8
\citep{blondel2008communities}. Hình~\ref{fig:communities} trình bày hồ sơ chủ
đề suy ra của mười hai cộng đồng lớn. Bài được loại trùng trong từng cộng đồng
và chia đóng góp giữa các nhãn. Một bài có thể gắn với nhiều cộng đồng qua các
tác giả của nó.

\plot{community-topic-profiles}{Hồ sơ chủ đề của mười hai cộng đồng đầu tiên
theo thứ tự kích thước đã lưu. Hình chọn mười chủ đề có tổng phần đóng góp lớn
nhất trong các hồ sơ này. Mỗi cộng đồng chỉ lưu năm chủ đề đứng đầu: ô xám là
không được báo cáo, không phải bằng không. Mã cộng đồng chỉ thuộc phân hoạch
cố định này.}{fig:communities}
\FloatBarrier

Hồ sơ cho thấy các hỗn hợp chủ đề khác nhau, không phải mỗi cộng đồng chỉ ứng
với một chủ đề. Học biểu diễn xuất hiện ở nhiều nhóm, trong khi một số chủ đề
khác tập trung hơn ở những hồ sơ nhất định. Sự giao thoa hữu ích khi tìm nhóm
nghiên cứu lân cận có mối quan tâm liên quan. Ranh giới cộng đồng phụ thuộc
đồ thị và tham số Louvain; đây là nhóm phân tích, không phải khoa hay ngành
chính thức, ổn định.

Độ trung tâm bổ sung một góc nhìn khác. Độ trung tâm bậc phản ánh phạm vi cộng
tác trực tiếp; độ mạnh phản ánh cộng tác lặp lại có trọng số phân bổ. Độ trung
tâm trung gian đo mức một nhà nghiên cứu nằm trên các đường đi ngắn nhất giữa
những người khác \citep{freeman1977betweenness}. Phân tích cố định dùng nghịch
đảo độ mạnh phân bổ dương làm độ dài cạnh và xấp xỉ độ trung tâm trung gian có
trọng số bằng 256 nút mốc. Bảng~\ref{tab:bridges} bổ sung ngữ cảnh cộng đồng
cho năm định danh đầu tiên, thay vì chỉ diễn giải thứ hạng.

% Generated from frozen v4; edit generator.
\begin{table}[!htbp]\centering\small
\caption{\lang{Leading approximate weighted-betweenness records and community context.}{Các định danh đứng đầu theo độ trung tâm trung gian xấp xỉ có trọng số và cộng đồng tương ứng.}}\label{tab:bridges}
\begin{tabularx}{\textwidth}{Yrrr}
\toprule \lang{Researcher}{Nhà nghiên cứu} & \lang{Betweenness}{\makecell[r]{Độ trung tâm\\trung gian}} & \lang{Community}{Cộng đồng} & \lang{Members}{Số thành viên}\\\midrule
Bo An & 0.042193 & 00126 & \num{367}\\
Lei Feng & 0.034010 & 00126 & \num{367}\\
Hau Chan & 0.021608 & 00000 & \num{687}\\
Shirui Pan & 0.020971 & 00043 & \num{481}\\
Chengjie Wang & 0.019882 & 00058 & \num{953}\\
\bottomrule\end{tabularx}\end{table}


Bo An (A5017743551) đứng đầu với 0.042193 trong 32,189 nhà nghiên cứu.
Định danh này thuộc cộng đồng 367 thành viên, gắn với 206 bài phân biệt, trong
đó 180 bài có nhãn. Học biểu diễn có 63 bài và 52.3333 đơn vị chủ đề phân bổ,
chiếm 29.1\% tổng phần đóng góp chủ đề của cộng đồng. Lei Feng, định danh đứng
thứ hai, thuộc cùng cộng đồng được chọn. Ví dụ này cho thấy nhiều định danh
cầu nối có thể cùng thuộc một bối cảnh hợp tác rộng hơn, thay vì là những cá
nhân dẫn đầu hoàn toàn tách biệt.

Hai bài có Bo An tham gia đã được kiểm tra nguồn là
\href{https://doi.org/10.1609/aaai.v36i8.20851}{\emph{GearNet: Stepwise Dual
Learning for Weakly Supervised Domain Adaptation}} (2022) và
\href{https://doi.org/10.1609/aaai.v37i6.25871}{\emph{Partial-Label Regression}}
(2023). Chúng minh họa sự tham gia được chấp nhận, nhưng riêng hai bài không
giải thích điểm trung gian.

\plot{researchers-network}{Bo An (màu vàng) và mười hai láng giềng một bước được
chấp nhận, chọn theo số bài hỗ trợ giảm dần, trọng số phân bổ và ID. Hiển thị
mọi cạnh được chấp nhận giữa mười ba ID đã chọn: 24 cạnh. Cạnh xanh lam chạm
Bo An và có độ rộng theo số bài hỗ trợ; cạnh xám nối các láng giềng. Hạt giống
bố trí là 17; hình không bao gồm toàn bộ mạng lân cận của Bo An.}{fig:ego}
\FloatBarrier

Mạng lân cận có cạnh giữa các đồng tác giả được chọn, không chỉ là hình sao
quanh Bo An. Vì vậy, hình hỗ trợ khảo sát bối cảnh hợp tác cục bộ; vai trò
trung gian toàn cục vẫn là thuộc tính của phép chiếu đầy đủ. Diễn giải còn phụ
thuộc quy tắc hỗ trợ: bộ lọc cố định chỉ giữ cặp có ít nhất hai bài giữ lại
7.78\% cạnh đồng tác giả và không có định danh nào chung với danh sách hai mươi
định danh đứng đầu độ trung tâm trung gian của mạng cơ sở. Bộ lọc biểu diễn
hợp tác lặp lại và làm mạng phân mảnh mạnh. Đây là lý do cần xem thứ hạng cầu
nối như gợi ý để khảo sát cấu trúc. Độ trung tâm trung gian đo đường đi theo
định nghĩa đồ thị đã chọn, không đo chất lượng nghiên cứu hay ảnh hưởng thực tế.

\section{Mạng trích dẫn trong tập dữ liệu}
Phép chiếu trích dẫn có hướng gồm toàn bộ \Papers{} bài và \CitationEdges{}
cạnh trích dẫn nội bộ được chấp nhận. Tham chiếu tới bài ngoài hai tập kỷ yếu
bị loại; ánh xạ đích mơ hồ và cạnh tự nối bị bỏ, còn tham chiếu lặp tới
cùng đích được gộp. Có \CitationIsolates{} bài không có cạnh trích dẫn nội bộ
vào hoặc ra; tổng bậc trung bình là 1.07 và bậc vào lớn nhất là 51.

\plot{top-cited-papers}{Mười bài có bậc vào cao nhất trong mạng trích dẫn thu
thập được. Mỗi số đếm là số bài trích dẫn phân biệt bên trong tập dữ liệu.
Giữ nguyên tiêu đề chính thức. Số đếm không gồm trích dẫn ngoài tập và chưa
chuẩn hóa theo tuổi công bố.}{fig:cited}
\FloatBarrier

\emph{Informer: Beyond Efficient Transformer for Long Sequence Time-Series
Forecasting} (AAAI 2021, \sourcepath{paper:aaai:2021:17325}) dẫn đầu với 51
bài trích dẫn nội bộ. Đây là ví dụ một bài dự báo chuỗi thời gian thu hút sự
chú ý lặp lại trong các kỷ yếu AI đã chọn. Thứ hạng giúp người đọc tìm điểm
xuất phát về phương pháp thường được tham chiếu, nhưng không phải bảng xếp
hạng trích dẫn toàn ngành.

Nhiều nút cô lập trong mạng trích dẫn cũng có ý nghĩa: giới hạn cả hai đầu
trong hai chuỗi hội nghị loại đi phần lớn mạng tài liệu tham khảo rộng hơn.
Một bài cô lập ở đây vẫn có thể trích dẫn và được trích dẫn nhiều ở nơi khác.
Bài cũ cũng có nhiều thời gian tích lũy trích dẫn hơn. Hai giới hạn này giúp
tránh diễn giải số trích dẫn nội bộ thấp thành mức chú ý tổng thể thấp.

\section{Các phát hiện chính}
\begin{enumerate}
\item \textbf{Quy mô công bố mở rộng nhưng không đều.} Số bài tăng từ 2,240
lên 4,042 sau bước giảm năm 2022. Tỷ trọng đa phương thức tăng trong phần có
nhãn ở cả hai hội nghị, tạo tín hiệu chủ đề đáng khảo sát cùng mức tăng số bài.
\item \textbf{Mạng đồng tác giả có lõi lớn cùng nhiều nhóm cục bộ.} Thành
phần lớn nhất chứa 20,776 trong 32,189 nhà nghiên cứu, trong khi toàn mạng có
2,936 thành phần. Kết quả hỗ trợ cả điều hướng mạng rộng và khảo sát từng nhóm.
\item \textbf{Cùng công bố liên khu vực có tính lặp lại.} Học thuật--doanh
nghiệp có 3,533 lượt hỗ trợ cặp--bài. Các ví dụ 14 bài về Học biểu diễn chỉ ra
những liên kết công bố cụ thể phục vụ lập kế hoạch nghiên cứu, trong phạm vi
bằng chứng tổ chức được giữ lại.
\item \textbf{Vị trí cầu nối và phạm vi cộng tác trả lời hai câu hỏi khác nhau.}
Bo An dẫn đầu độ trung tâm trung gian xấp xỉ với 0.042193, nhưng bộ lọc hỗ trợ
lặp lại thay đổi toàn bộ hai mươi vị trí đầu. Danh sách cần đi kèm phép chiếu
và quy tắc hỗ trợ, không chỉ tên người.
\item \textbf{Sự chú ý qua trích dẫn nằm trong một cửa sổ hẹp.} Informer có
bậc vào lớn nhất là 51, trong khi 7,302 bài là nút cô lập trích dẫn. Cả hai
kết quả phản ánh giới hạn hai hội nghị và danh sách tham chiếu thu thập được.
\end{enumerate}
